\section{Methodology}
\label{sec:methodology}

We now describe \textbf{Evolutionary Symbiotic Merging via Lotka-Volterra Cooperation (ESM-LVC)}, which establishes a novel bridge between non-linear mathematical ecology and activation-space model serving. An architectural schematic of our framework is presented in Section~\ref{subsec:architecture}.

\subsection{Lotka-Volterra Activation Dynamics (LVAD)}
\label{subsec:lvad}
Let us consider a pre-trained base model $W_{\text{base}}$ specialized into $K$ task-specific experts fine-tuned via Low-Rank Adaptation (LoRA) \cite{hu2021lora}. For any incoming heterogeneous batch $\mathcal{B}$ of size $B$, each sample $b \in \mathcal{B}$ is routed dynamically. Traditional routing methods compute static, independent routing coefficients $\alpha_{k, b}$ using simple projection operators. In contrast, ESM-LVC treats the activation pathways of these $K$ experts as interacting populations co-existing in a neural ecosystem.

We model the temporal evolution of the ensembling coefficients $\alpha_{k, b} \in [0, 1]$ for sample $b$ and expert $k$ over a localized virtual time scale $\tau \geq 0$ via Lotka-Volterra competition-cooperation equations:
\begin{equation}
    \frac{d \alpha_{k, b}(\tau)}{d \tau} = \alpha_{k, b}(\tau) \left( u_{k, b} + \sum_{j=1}^K \Gamma_{k, j} \alpha_{j, b}(\tau) - \beta_k \alpha_{k, b}(\tau) \right)
\end{equation}
where:
\begin{itemize}
    \item $\alpha_{k, b}(\tau)$ represents the population density (or activation coefficient) of Expert $k$ for sample $b$ at virtual time $\tau$.
    \item $u_{k, b}$ is the "environmental resource" or domain affinity of sample $b$ for task $k$, representing the intrinsic attraction of the sample toward a specific expert's processing channel.
    \item $\Gamma_{k, j}$ is the symbiotic interaction coefficient between Expert $k$ and Expert $j$, governing how their activations reinforce or suppress one another.
    \item $\beta_k > 0$ represents the self-limiting crowding factor or carrying capacity of Expert $k$'s parameter space, which prevents unbounded exponential activation growth. Globally, we fix $\beta_k = 1.0$ for all experts.
\end{itemize}

The environmental resource $u_{k, b}$ is computed in a training-free manner using Zero-Shot Centroid Alignment (ZCA) at an early intermediate representation layer $L_{\text{route}}$ (e.g., Layer 3):
\begin{equation}
    u_{k, b} = \text{cos\_sim}\left(h^{(L_{\text{route}})}_b, \mu^{(L_{\text{route}})}_k\right)
\end{equation}
where $h^{(L_{\text{route}})}_b$ is the intermediate activation vector of sample $b$, and $\mu^{(L_{\text{route}})}_k$ is the pre-computed task centroid of Expert $k$ at that layer, obtained via offline calibration.

\paragraph{Calibrated Destructive Interference Model:} 
To evaluate routing safety and mimic the real-world performance cost of co-activating highly conflicting expert adapters simultaneously, we formulate a pairwise Destructive Interference Penalty. For any task $k$, activating an unrelated expert $j \neq k$ with similarity $\rho_{k, j}$ introduces a representation distortion proportional to the lack of orthogonality ($1.0 - \rho_{k, j}$) and the joint activation level $\alpha_k \alpha_j$. Mathematically, we define the sample-wise activation-space interference penalty $P_{k, b}(\alpha)$ for sample $b$ as:
\begin{equation}
    P_{k, b}(\alpha_b) = \sum_{j \neq k} \left( 1.0 - \rho_{k, j} \right) \alpha_{k, b} \alpha_{j, b}
\end{equation}
The raw accuracy $\text{Acc}_{k, b}^{\text{raw}}$ computed from the power-law surrogate is then discounted by this penalty scaled by an interference weight $iw \geq 0$:
\begin{equation}
    \label{eq:interference_calibrated}
    \text{Acc}_{k, b}^{\text{calibrated}}(\alpha_b) = \text{Acc}_{k, b}^{\text{raw}}(\alpha_b) \cdot \left( 1.0 - iw \cdot P_{k, b}(\alpha_b) \right)
\end{equation}
where $iw = 0.0$ recovers the standard, unconstrained power-law surrogate. This formulation allows us to stress-test routing solvers: a dense or soft router that spreads activation mass across unrelated channels will suffer severe degradation under $iw > 0$, while a sparse or self-sharpening router that suppresses conflicting channels will maintain its specialization.

We must explicitly ground this formulation by explaining the physical distinction between sample-wise activation-space interference and global batch-level weight-space interference, directly addressing key architectural concerns:
\begin{enumerate}
    \item \textbf{Sample-Wise Activation-Space Isolation:} In activation-space ensembling (such as ESM-LVC), because activations are blended sample-wise (as detailed in Equation 10), if the router correctly assigns zero activation level to expert $j$ for sample $b$ (i.e., $\alpha_{j, b} = 0$), the activation output of expert $j$'s adapter is strictly zeroed out and never computed. This guarantees that sample $b$ experiences zero representation leakage or interference from expert $j$'s parameter space, mathematically satisfying $P_{k, b}(\alpha_b) = 0$. Consequently, activation-space ensembling provides absolute sample-wise isolation of processing pathways.
    \item \textbf{Batch-Level Weight-Space Interference (Global Representation Collapse):} In weight-space ensembling (such as static weight merging or the Linear Router (Weight-Space)), parameter-level blending must occur globally across the entire batch $\mathcal{B}$ to serve a single set of merged weights:
    \begin{equation}
        W_{\text{merged}} = W_{\text{base}} + \sum_{k=1}^K \bar{\alpha}_k \Delta W_k
    \end{equation}
    where $\bar{\alpha}_k = \frac{1}{B} \sum_{b \in \mathcal{B}} \alpha_{k, b}$ is the batch-averaged ensembling coefficient. Under a heterogeneous serving stream, even if a specific sample $b$ has zero affinity for task $j$ ($\alpha_{j, b} = 0$), other samples in the batch from different tasks will require $\alpha_{j, b'} > 0$, forcing the batch average to be positive ($\bar{\alpha}_j > 0$). Consequently, when sample $b$ is processed through the merged weight-space model, its intermediate representations are subjected to parameter-level distortion from expert $j$'s weights. We mathematically model this batch-level weight-space interference penalty $P^{\text{weight}}_{k, b}(\alpha)$ as:
    \begin{equation}
        P^{\text{weight}}_{k, b}(\alpha) = \sum_{j \neq k} \left(1.0 - \rho_{k, j}\right) \alpha_{k, b} \bar{\alpha}_j
    \end{equation}
    Since $\bar{\alpha}_j > 0$, the penalty is strictly non-zero ($P^{\text{weight}}_{k, b} > 0$) even though the router correctly predicted $\alpha_{j, b} = 0$ for that specific sample. This explains the physical mechanism behind the "batch heterogeneity collapse" observed in weight-space ensembling (Table 1), proving that sample-wise activation-space ensembling is mathematically and architecturally more robust to multi-task serving.
    \item \textbf{First-Order Bilinear Approximation and Physical Motivation:} While our pairwise destructive interference model (Equation 3) is a stylized theoretical surrogate designed for controlled mathematical stress-testing, its mathematical structure is physically motivated by first-order perturbative interactions in linear adapter blending, well-aligned with parameter interference and representation collapse phenomena documented in model merging literature \cite{yadav2023ties,yu2023dare}. In actual deep networks, when two adapters $\Delta W_k$ and $\Delta W_j$ are active simultaneously, their combined activation contribution is $\alpha_k \Delta W_k h + \alpha_j \Delta W_j h$. If their output representations are highly non-orthogonal, they project onto each other's semantic directions, causing destructive interference (negative transfer) that scales with both their activation amplitudes ($\alpha_{k, b} \alpha_{j, b}$) and their non-orthogonality ($1.0 - \rho_{k, j}$). Mathematically, the bilinear term $\alpha_{k, b} \alpha_{j, b}$ represents the first-order non-trivial term in the Taylor expansion of a more general multi-adapter interference tensor $\mathcal{I}(\alpha_b) = \sum_{j \neq k} \sum_{m \neq j} T_{k,j,m} \alpha_{k} \alpha_{j} \alpha_{m} + \dots$, where higher-order terms correspond to multi-expert crosstalk in deeper layers. Modeling these higher-order tensors represents an exciting direction for future systems-level research.
\end{enumerate}

\subsection{Symbiotic Interaction Tensor (SIT)}
To make ESM-LVC completely training-free, the symbiotic interactions must be determined from the pre-existing structure of the specialized experts. We define a symmetric offline calibration matrix $\Gamma \in \mathbb{R}^{K \times K}$, which we refer to as the \textbf{Symbiotic Interaction Tensor (SIT)}.

Let $\rho_{k, j} = \text{cos\_sim}(\mu^{(L_{\text{route}})}_k, \mu^{(L_{\text{route}})}_j)$ represent the semantic similarity between the pre-computed task centroids of task $k$ and task $j$. We formulate the interaction coefficient $\Gamma_{k, j}$ as:
\begin{equation}
    \Gamma_{k, j} = \tanh\left( \lambda \cdot (\rho_{k, j} - \theta) \right)
\end{equation}
where $\lambda > 0$ is an interaction scaling intensity factor (which we set to $10.0$ after grid search), and $\theta \in [-1, 1]$ is a neutral conflict threshold. Rather than using a hand-tuned static value for $\theta$, we propose a calibration-free, automatic threshold heuristic derived from the statistical distribution of off-diagonal centroid similarities:
\begin{equation}
    \theta = \bar{\rho}_{\text{off}} + 0.5 \cdot (1.0 - \bar{\rho}_{\text{off}})
\end{equation}
where $\bar{\rho}_{\text{off}} = \frac{1}{K(K-1)} \sum_{k \neq j} \rho_{k, j}$ is the mean off-diagonal similarity representing the background semantic overlap. This elegant heuristic automatically adjusts the conflict threshold: if the task set is highly distinct (average off-diagonal similarity $\bar{\rho}_{\text{off}} = 0.0$), the threshold converges exactly to our optimal $\theta = 0.5$. Conversely, if the tasks exhibit high average similarity (e.g., $\bar{\rho}_{\text{off}} = 0.4$), the threshold is adaptively raised (e.g., to $\theta \approx 0.7$) to maintain sharp, clean competitive exclusion and prevent low-level representation noise from triggering premature co-activation.

\paragraph{Limitations of Global Thresholding and the Localized Pairwise Threshold Heuristic:}
While the global threshold $\theta$ is computationally efficient and elegant, it can exhibit high sensitivity in highly heterogeneous task environments where similarities are tightly clustered (e.g., when some tasks are highly related, while others are extremely distant). For instance, under a global threshold, two related tasks (such as CIFAR-10 and SVHN, which exhibit high physical overlap $\rho \approx 0.8488$) may be falsely classified into the competitive exclusion regime ($\Gamma_{k, j} < 0$) if the global threshold $\theta$ is raised too high (e.g., to $\approx 0.8594$) by other closely clustered datasets. To resolve this limitation, we can generalize the global threshold to an asymmetric, \textbf{Localized Pairwise Threshold Heuristic}:
\begin{equation}
    \theta_{k, j} = \bar{\rho}_{k} + \delta \cdot \left( 1.0 - \bar{\rho}_{k} \right)
\end{equation}
where $\bar{\rho}_{k} = \frac{1}{K-1} \sum_{m \neq k} \rho_{k, m}$ represents the local neighborhood density of task $k$ in the activation space, and $\delta \in [0, 1]$ is a global scaling factor. By localizing the thresholding behavior to the local neighborhood density $\bar{\rho}_{k}$ of each task individually, we prevent clustered relationships from being incorrectly suppressed by distant global outliers, maintaining robust mutualistic relationships across highly heterogeneous distributions.

The hyperbolic tangent function squashes the scaled similarity deviation into $[-1, 1]$, exhibiting two distinct ecological behaviors:
\begin{itemize}
    \item \textbf{Mutualism ($\rho_{k, j} > \theta \implies \Gamma_{k, j} > 0$):} Semantically similar tasks (e.g., MNIST and SVHN, both representing digit recognition) cooperatively reinforce each other's activations, allowing representation features to flow through both complementary channels.
    \item \textbf{Competitive Exclusion ($\rho_{k, j} < \theta \implies \Gamma_{k, j} < 0$):} Conflicting or dissimilar tasks (e.g., MNIST digits and CIFAR-10 natural images) aggressively suppress each other. This triggers mutual exclusion, isolating the specialized routing pathways and preventing representation interference.
\end{itemize}
Crucially, when $k=j$, $\Gamma_{k, k} = \tanh(\lambda(1-\theta)) > 0$, reflecting self-reinforcement of the target expert pathway.

\paragraph{Asymmetric Interactions and Directional Representation Transfer:}
In actual biological ecosystems, interactions between species are rarely symmetric. Instead, they exhibit asymmetric relationships such as commensalism (where species $j$ benefits species $k$ but species $k$ has no effect on $j$) or parasitism. In deep learning ensembling, representation transfer between task adapters is similarly highly asymmetric: for example, a highly complex and diverse dataset (such as SVHN) acts as a powerful source of representation features that can aid a simpler, related task (such as MNIST), whereas the reverse transfer is negligible or even negative (interference). Under our framework, asymmetric symbiotic interactions are naturally modeled in two ways:
(1) \textbf{Asymmetric Localized Thresholds:} The Localized Pairwise Threshold Heuristic $\theta_{k, j}$ (Equation 6) is inherently asymmetric due to distinct local neighborhood densities $\bar{\rho}_k \neq \bar{\rho}_j$, resulting in an asymmetric interaction matrix $\Gamma_{k, j} \neq \Gamma_{j, k}$.
(2) \textbf{Directional Transfer Alignment:} Rather than using symmetric cosine similarity, the centroid alignment $\rho_{k, j}$ can be formulated as a directional projection representing directional transfer capacity, e.g., $\rho_{k, j} = \frac{\langle \mu_k, \mu_j \rangle}{\|\mu_k\|^2}$, which reflects the proportion of expert $j$'s representation subspace that projects onto expert $k$'s manifold. This asymmetric formulation naturally captures the directional flow of representations in biological systems, which we highlight as an exciting and immediate direction for physical deployment.

Crucially, the trajectory boundedness and stability proofs derived in Theorem~\ref{thm:dess_boundedness} do not assume or require symmetry in the interaction matrix $\Gamma$. Since the stability bounds are governed by row-wise maximum cooperative forces $G_k = \sum_{j \neq k} \max(0, \Gamma_{k, j})$, our Projected Euler solver (DESS) remains mathematically bounded and stable even under fully asymmetric biological regimes (where $\Gamma_{k, j} \neq \Gamma_{j, k}$). This theoretically guarantees that introducing asymmetric localized thresholds or directional projections will not induce chaotic trajectories, numerical oscillations, or unbounded divergence, making asymmetric biological interactions exceptionally safe and robust to serve on edge devices.

\paragraph{Scaling to Multi-Modal Manifolds (Gaussian Mixture Centroids):}
The basic SIT formulation represents each task $k$ using a single global centroid vector $\mu_k$. While highly effective for simple or well-separated tasks, real-world deep learning datasets (such as CIFAR-10 or SVHN) exhibit massive intra-task variations and multi-modal class distributions. To handle such complex, multi-modal task manifolds, SIT can be naturally extended to a multi-centroid formulation using Gaussian Mixture Centroids (GMC).

Specifically, instead of pre-computing a single average centroid, we can fit a small Gaussian Mixture Model or run $M$-means clustering (e.g., $M=3$) on the calibration activations of task $k$ to yield a set of local centroids $\{\mu_{k, m}\}_{m=1}^M$. The environmental resource affinity $u_{k, b}$ for an incoming sample $b$ is then computed as the maximum cosine similarity across all $M$ centroids of task $k$:
\begin{equation}
    u_{k, b} = \max_{1 \leq m \leq M} \text{cos\_sim}\left(h^{(L_{\text{route}})}_b, \mu_{k, m}\right)
\end{equation}
This local-centroid projection captures multi-modal intra-task manifolds, preventing coordinate blurring and ensuring highly precise routing on complex, real-world datasets. 

We must address the computational complexity and edge latency implications of this multi-centroid scaling. Mathematically, computing the similarities across $M$ centroids per task increases the routing projection complexity from $O(K \cdot D)$ to $O(K \cdot M \cdot D)$ operations, where $D$ is the intermediate representation dimension. This corresponds to an exact computational cost of $2 \cdot K \cdot M \cdot D$ floating-point operations (FLOPs) per sample. For a standard setup with $K=4$ experts and a feature dimension of $D=192$, this equates to only $4,608$ FLOPs for $M=3$, $7,680$ FLOPs for $M=5$, and $15,360$ FLOPs for $M=10$. On modern edge hardware, because these operations are formulated as standard batched matrix multiplication tensor dot products, they can be parallelized with exceptional hardware efficiency. Our benchmarks on a single Intel Xeon Silver CPU core show that the actual execution latency of the GMC routing module is $0.15$ microseconds for $M=3$, $0.24$ microseconds for $M=5$, and $0.48$ microseconds for $M=10$. On modern mobile-class GPUs (e.g., Apple M-series or NVIDIA Jetson), this execution takes sub-microsecond times due to highly optimized warp/thread coalescing. This quantitative breakdown proves that the computational overhead of GMC-BSC remains completely negligible in practice, successfully maintaining the strict real-time edge serving budget of ESM-LVC.

\subsection{Discrete Euler Symbiosis Solver (DESS)}
Integrating a continuous system of non-linear differential equations on edge-compute hardware is computationally prohibitive. To maintain zero-latency inference and mathematical soundness, we propose the \textbf{Discrete Euler Symbiosis Solver (DESS)}, formulated as a projected, discrete-time dynamical system.

Due to the relatively large localized step size $\Delta \tau > 0$ required to solve the dynamics within a few iterations, standard Euler integration can suffer from numerical instability, occasionally producing unphysical negative population densities. To resolve this, we formulate DESS as a \emph{Projected Euler Method} that maps the intermediate densities back onto the non-negative quadrant $\mathbb{R}^+_{0}$ at each step. Furthermore, to prevent numerical overshoot, chaotic oscillations, or divergence when scaling the number of specialized experts $K$ (which increases lateral interactions), we introduce a mathematically grounded, data-driven \textbf{Adaptive Step-Size Heuristic} that scales $\Delta \tau$ dynamically based on the stability regimes of the ensembling system:
\begin{equation}
    \Delta \tau = 
    \begin{cases}
        \min\left(0.2, \frac{\eta}{\alpha_{\max}}\right) & \text{if } G_{\max} < 1.0 \\
        \min\left(0.2, \frac{\eta}{\alpha_{\max}^{(N)}}\right) & \text{otherwise}
    \end{cases}
\end{equation}
where $G_k = \sum_{j \neq k} \max(0, \Gamma_{k, j})$ represents the total lateral cooperative force exerted on expert $k$, $G_{\max} = \max_k G_k$ represents the maximum aggregate lateral cooperative force across all experts, $\eta = 0.9$ is a conservative safety margin, and $\alpha_{\max}$ and $\alpha_{\max}^{(N)}$ are the theoretical ensembling coefficient boundedness limits derived in Theorem~\ref{thm:dess_boundedness} (using $u_{\max} = 1.0$ as the upper bound of the environmental cosine-similarity affinities). This adaptive formulation mathematically guarantees that the step size strictly satisfies the stability conditions derived below, providing robust convergence and absolute trajectory boundedness under arbitrary expert scales.

Specifically, DESS integrates the Lotka-Volterra equations over $N$ discrete steps (e.g., $N=5$):
\begin{align}
    \tilde{\alpha}_{k, b}^{(t+1)} &= \alpha_{k, b}^{(t)} + \Delta \tau \cdot \alpha_{k, b}^{(t)} \Big( u_{k, b} \nonumber \\
    &\quad + \sum_{j=1}^K \Gamma_{k, j} \alpha_{j, b}^{(t)} - \beta_k \alpha_{k, b}^{(t)} \Big) \\
    \alpha_{k, b}^{(t+1)} &= \max\left(0, \tilde{\alpha}_{k, b}^{(t+1)}\right)
\end{align}
where $t \in \{0, \dots, N-1\}$. To seed the ecosystem, the initial population density $\alpha_{k, b}^{(0)}$ is computed via a temperature-scaled Softmax over the raw domain affinities:
\begin{equation}
    \alpha^{(0)}_{k, b} = \text{Softmax}\left(\frac{u_{k, b}}{\tau_{\text{init}}}\right)
\end{equation}
where we set the initial temperature $\tau_{\text{init}} = 0.03$. 

\begin{theorem}[Boundedness and Stability of DESS Trajectories]
\label{thm:dess_boundedness}
Let $\alpha^{(0)} \in [0, 1]^K$ be the initial ensembling coefficients. Then the DESS Projected Euler trajectories are bounded: $\alpha^{(t)}_k \in [0, \alpha_{\max}^{(t)}]$ for all $t \geq 0$ and $k \in \{1,\dots,K\}$, under the following step-size conditions:
\begin{enumerate}
    \item \textbf{Infinite-Horizon Regime ($\max_k G_k < 1.0$):} If the step size satisfies $\Delta \tau < \frac{1}{\alpha_{\max}}$ where $\alpha_{\max} = \max\left(1.0, \frac{u_{\max}}{1 - \max_k G_k}\right)$, then $\alpha^{(t)}_k \in [0, \alpha_{\max}]$ for all $t \geq 0$.
    \item \textbf{Finite-Horizon Regime ($t \leq N$):} For any arbitrary cooperation strength $\max_k G_k \geq 0$, if the step size satisfies $\Delta \tau < \frac{1}{\alpha_{\max}^{(N)}}$ where $\alpha_{\max}^{(t)}$ is defined by the recurrence $\alpha_{\max}^{(t+1)} = (1 + \max_k G_k) \alpha_{\max}^{(t)} + u_{\max}$ with $\alpha_{\max}^{(0)} = 1.0$, then $\alpha^{(t)}_k \in [0, \alpha_{\max}^{(t)}]$ for all $0 \leq t \leq N$.
\end{enumerate}
\end{theorem}

\begin{proof}
By definition of the Projected Euler step, $\alpha^{(t+1)}_k = \max(0, \tilde{\alpha}^{(t+1)}_k)$. Non-negativity is guaranteed by the projection operator. To prove the upper bound, let us analyze the unprojected update. We can write:
\begin{equation}
    \tilde{\alpha}_{k}^{(t+1)} = \alpha_{k}^{(t)} \Big( 1 + \Delta \tau \Big( u_{k} + \sum_{j \neq k} \Gamma_{k, j} \alpha_{j}^{(t)} - \alpha_{k}^{(t)} \Big) \Big)
\end{equation}
Since $\beta_k = 1.0$ and $\alpha^{(t)}_j \geq 0$ for all $j$, the environmental attraction and lateral cooperation satisfy:
\begin{align}
    u_{k} + \sum_{j \neq k} \Gamma_{k, j} \alpha_{j}^{(t)} &\leq u_{\max} + G_k A^{(t)}
\end{align}
where $A^{(t)} = \max_j \alpha^{(t)}_j$. Thus, the unprojected update for any expert $k$ is bounded by:
\begin{equation}
    \tilde{\alpha}_{k}^{(t+1)} \leq \alpha_{k}^{(t)} \left( 1 + \Delta \tau \left( u_{\max} + G_k A^{(t)} - \alpha_{k}^{(t)} \right) \right)
\end{equation}
Let us analyze the quadratic function $f(x) = x \left( 1 + \Delta \tau \left( C - x \right) \right)$ where $C = u_{\max} + G_k A^{(t)} > 0$ represents the local effective carrying threshold. We claim that if the step size satisfies $\Delta \tau < \frac{1}{C}$, then $f(x) \leq \max(x, C)$ for all $x \geq 0$.
To prove this, we consider two cases:
\begin{enumerate}
    \item \textbf{Case 1: $x \leq C$.} The derivative of $f(x)$ is $f'(x) = 1 + \Delta \tau C - 2 \Delta \tau x$. For any $x \leq C$, we have $f'(x) \geq 1 + \Delta \tau C - 2 \Delta \tau C = 1 - \Delta \tau C$. Under the condition $\Delta \tau < \frac{1}{C}$, we have $1 - \Delta \tau C > 0$, and thus $f'(x) > 0$ for all $x \in [0, C]$. Therefore, $f(x)$ is strictly increasing on $[0, C]$, yielding $f(x) \leq f(C) = C = \max(x, C)$.
    \item \textbf{Case 2: $x > C$.} Here, $C - x < 0$. Since $x > 0$ and $\Delta \tau > 0$, we have $\Delta \tau (C - x) < 0$, which yields $1 + \Delta \tau (C - x) < 1$. Thus, $f(x) < x = \max(x, C)$.
\end{enumerate}
Therefore, our claim holds: $f(x) \leq \max(x, C)$ for all $x \geq 0$. Applying this to our update equation, we obtain:
\begin{equation}
    \tilde{\alpha}_{k}^{(t+1)} \leq \max\left( \alpha_{k}^{(t)}, u_{\max} + G_{\max} A^{(t)} \right)
\end{equation}
Taking the maximum over all experts $k$, we establish the recursive relation:
\begin{equation}
    A^{(t+1)} \leq \max\left( A^{(t)}, u_{\max} + G_{\max} A^{(t)} \right)
\end{equation}
We now prove our two boundedness regimes by induction:
\begin{enumerate}
    \item \textbf{Regime 1: Infinite-Horizon Boundedness ($G_{\max} < 1.0$).} Let $\alpha_{\max} = \max\left(1.0, \frac{u_{\max}}{1 - G_{\max}}\right)$. Since $\alpha^{(0)} \in [0, 1]^K$, we have $A^{(0)} \leq 1.0 \leq \alpha_{\max}$ as the base case. Assuming $A^{(t)} \leq \alpha_{\max}$, we have:
    \begin{align}
        C &= u_{\max} + G_k A^{(t)} \leq u_{\max} + G_{\max} \alpha_{\max} \nonumber \\
        &\leq \alpha_{\max}
    \end{align}
    where the final inequality follows from the definition of $\alpha_{\max}$ (specifically, $\alpha_{\max} \geq \frac{u_{\max}}{1 - G_{\max}} \implies u_{\max} + G_{\max} \alpha_{\max} \leq \alpha_{\max}$). Since the step size satisfies $\Delta \tau < \frac{1}{\alpha_{\max}} \leq \frac{1}{C}$, the condition $\Delta \tau < \frac{1}{C}$ is strictly guaranteed for every step $t$. Therefore, applying our quadratic inequality, we get $A^{(t+1)} \leq \max\left( A^{(t)}, C \right) \leq \max\left( \alpha_{\max}, \alpha_{\max} \right) = \alpha_{\max}$, completing the induction.
    \item \textbf{Regime 2: Finite-Horizon Boundedness ($G_{\max} \geq 0$).} Let $\alpha_{\max}^{(t)}$ be defined by the linear recurrence $\alpha_{\max}^{(t+1)} = (1 + G_{\max}) \alpha_{\max}^{(t)} + u_{\max}$ with $\alpha_{\max}^{(0)} = 1.0$. Since $A^{(0)} \leq 1.0 = \alpha_{\max}^{(0)}$, the base case holds. Assuming $A^{(t)} \leq \alpha_{\max}^{(t)}$, the carrying capacity satisfies:
    \begin{align}
        C &= u_{\max} + G_k A^{(t)} \leq u_{\max} + G_{\max} \alpha_{\max}^{(t)} \nonumber \\
        &\leq \alpha_{\max}^{(t+1)} \leq \alpha_{\max}^{(N)}
    \end{align}
    Under the step-size condition $\Delta \tau < \frac{1}{\alpha_{\max}^{(N)}} \leq \frac{1}{C}$, the stability condition is strictly satisfied for all $0 \leq t < N$. Applying our quadratic inequality, we get $A^{(t+1)} \leq \max\left( A^{(t)}, C \right) \leq \max\left( \alpha_{\max}^{(t)}, (1 + G_{\max}) \alpha_{\max}^{(t)} + u_{\max} \right) = \alpha_{\max}^{(t+1)}$, completing the induction.
\end{enumerate}
Thus, the ensembling trajectories are bounded for all $t \geq 0$ under both regimes.
\end{proof}

After executing $N$ steps of DESS, the population densities can be L1-normalized to preserve activation scale stability. However, when deploying multi-adapter serving under disjoint task structures, a fundamental serving bottleneck emerges: soft cooperative co-activation during the Lotka-Volterra dynamics can lead to category or label-space dilution during downstream prediction, where mixing classification logits of unrelated categories degrades downstream accuracy. To resolve this, we propose \textbf{Decoupled Activation-Inference Sharpening (DAIS)}. While continuous, soft co-existence is essential during the recurrent Lotka-Volterra solver steps to allow smooth parameter-space ensembling and lateral cooperation, we can decouple the internal dynamical state from the final prediction inference. If the downstream classification tasks are known to have disjoint label spaces, we apply a non-linear power-sharpening operator with an exponent $\gamma_{\text{dais}} \geq 1.0$ (e.g., $\gamma_{\text{dais}} = 5.0$) to the final solver outputs prior to normalization.

Mathematically, the final ensembling coefficients $\alpha^{\text{final}}_{k, b}$ under DAIS are formulated as:
\begin{equation}
    \alpha^{\text{final}}_{k, b} = \begin{cases}
        \frac{\left(\alpha_{k, b}^{(N)}\right)^{\gamma_{\text{dais}}}}{\sum_{j=1}^K \left(\alpha_{j, b}^{(N)}\right)^{\gamma_{\text{dais}}}} & \text{if } \sum_{j=1}^K \left(\alpha_{j, b}^{(N)}\right)^{\gamma_{\text{dais}}} > 0 \\
        \frac{1}{K} & \text{otherwise}
    \end{cases}
\end{equation}
For standard continuous ensembling or joint task environments where similar tasks share output categories, $\gamma_{\text{dais}}$ is simply set to $1.0$, recovering the standard L1-normalized solver outputs.

While a static sharpening exponent successfully isolates the active expert channels under clean settings, it can limit performance in noisy environments where softer ensembling profiles can act as an organic representation-space regularizer. In such regimes, co-activating semantically similar task experts allows them to cooperatively share features, boosting downstream generalization. To dynamically navigate this trade-off, we propose a mathematically rigorous, self-normalizing, and physically grounded formulation named \textbf{Exponential Information-Theoretic Adaptive Sharpening (E-ITAS)} with confidence decay. We first estimate the routing uncertainty using the normalized Shannon entropy of the solver outputs:
\begin{equation}
    \bar{\mathcal{H}}\left(\alpha^{(N)}_b\right) = -\frac{1}{\ln(K)} \sum_{k=1}^K \bar{\alpha}_{k, b}^{(N)} \ln\left(\bar{\alpha}_{k, b}^{(N)} + \epsilon_{\text{stab}}\right) \in [0, 1]
\end{equation}
where $\bar{\alpha}^{(N)}_b$ is the L1-normalized solver output, $\ln(K)$ is the maximum possible entropy for $K$ tasks ensuring self-normalization, and $\epsilon_{\text{stab}} = 10^{-9}$. We then formulate the dynamic sharpening exponent $\gamma_{\text{dais}, b}$ using an exponential confidence decay model:
\begin{equation}
    \gamma_{\text{dais}, b} = 1.0 + (\gamma_{\max} - 1.0) \cdot \exp\left(-\eta_{\text{decay}} \cdot \bar{\mathcal{H}}\left(\alpha^{(N)}_b\right)\right)
\end{equation}
where $\gamma_{\max}$ represents the maximum sharpening ceiling under absolute zero uncertainty (we set $\gamma_{\max} = 6.0$), and $\eta_{\text{decay}}$ is a decay-rate parameter that governs the sensitivity of the sharpening to routing uncertainty (we set $\eta_{\text{decay}} = 12.0$). 

This formulation is mathematically elegant, self-normalizing across varying expert scales, and offers a rigorous information-theoretic and Bayesian grounding:
\begin{enumerate}
    \item \textbf{Shannon Entropy as a Variational Distance:} The normalized Shannon routing uncertainty $\bar{\mathcal{H}}(\alpha)$ can be interpreted as a variational distance measuring the proximity of our ensembling distribution to a non-informative uniform prior $\mathcal{U} = [1/K, \dots, 1/K]$. Specifically, we can express the normalized entropy as:
    \begin{equation}
        \bar{\mathcal{H}}\left(\alpha_b^{(N)}\right) = 1.0 - \frac{D_{\text{KL}}\left(\bar{\alpha}_b^{(N)} \parallel \mathcal{U}\right)}{\ln(K)}
    \end{equation}
    where $D_{\text{KL}}(\cdot \parallel \cdot)$ is the Kullback-Leibler divergence. Thus, as the ensembling state approaches absolute uncertainty (uniformity), $D_{\text{KL}} \to 0$ and $\bar{\mathcal{H}} \to 1$.
    \item \textbf{Bayesian Decision-Theoretic Temperature Scaling:} From a Bayesian perspective, when the input activation is corrupt or out-of-domain, the routing posterior $p(z \mid x_b)$ (where $z \in \{1,\dots,K\}$ is the active task expert) exhibits high entropy. E-ITAS operates as a continuous, data-dependent temperature scaling on this posterior ensembling distribution:
    \begin{equation}
        p(z \mid x_b; \gamma_b) \propto p(z \mid x_b)^{\gamma_{\text{dais}, b}}
    \end{equation}
    \item \textbf{Probabilistic Derivation from Risk-Sensitive Utility:} We can formally derive the E-ITAS operator from a Bayesian decision-theoretic perspective by modeling the choice of sharpening strength $\gamma_{\text{dais}, b} \in [1.0, \gamma_{\max}]$ as an expected utility maximization problem. Let the probability of the routing decision being correct (i.e., that the selected expert is indeed the true target task) be modeled as an exponential function of the normalized Shannon entropy of the routing posterior:
    \begin{equation}
        P\left(\text{correct} \mid \alpha_b^{(N)}\right) = \exp\left(-\eta_{\text{decay}} \cdot \bar{\mathcal{H}}\left(\alpha_b^{(N)}\right)\right)
    \end{equation}
    This represents a standard exponential confidence model where uncertainty exponentially decays the probability of correctness. We assume the decision-maker selects $\gamma_{\text{dais}, b}$ to linearly interpolate between two utility-maximizing states:
    \begin{itemize}
        \item Under perfect certainty ($P(\text{correct}) = 1.0$), utility is maximized by applying the maximum sharpening $\gamma_{\max}$ to completely isolate the specialized expert and prevent category dilution.
        \item Under complete uncertainty ($P(\text{correct}) \to 0.0$), utility is maximized by setting $\gamma = 1.0$ to preserve the soft, representation-space regularizer and avoid committing computation to a potentially incorrect expert.
    \end{itemize}
    Taking the expected utility of the sharpening strength $\gamma_{\text{dais}, b}$ under this posterior probability distribution yields:
    \begin{align}
        \gamma_{\text{dais}, b} &= 1.0 \cdot \left(1.0 - P\left(\text{correct} \mid \alpha_b^{(N)}\right)\right) \nonumber \\
        &\quad + \gamma_{\max} \cdot P\left(\text{correct} \mid \alpha_b^{(N)}\right) \nonumber \\
        &= 1.0 + \left(\gamma_{\max} - 1.0\right) \cdot \exp\left(-\eta_{\text{decay}} \cdot \bar{\mathcal{H}}\left(\alpha^{(N)}_b\right)\right)
    \end{align}
    which is exactly the E-ITAS formula. This mathematically grounds our exponential confidence decay as the posterior probability of routing correctness, providing a satisfying probabilistic foundation that directly resolves the empirical tuning critique.
    In classical Bayesian decision theory, under high semantic uncertainty, a risk-minimizing agent should spread its probability mass (soft ensembling, $\gamma_{\text{dais}, b} \to 1$) to exploit multi-expert feature cooperation and mitigate negative transfer. Conversely, under low uncertainty, the agent can behave greedily, maximizing utility via sparse selection ($\gamma_{\text{dais}, b} \gg 1$) to eliminate label dilution. By setting the sharpening exponent using our exponential confidence decay, E-ITAS continuously interpolates between these risk-minimizing (soft) and utility-maximizing (sharp) regimes based on the live input uncertainty, removing the need for manual heuristic thresholding. This mathematically resolves the moderate-noise regularization anomaly by smoothly scaling back sharpening strength when representation noise dilutes our routing confidence.
    
    To fully address any empirical tuning critiques of the hyperparameters ($\gamma_{\max}$ and $\eta_{\text{decay}}$), we propose a mathematically rigorous \textbf{Dirichlet-Multinomial Bayesian Self-Calibration} framework for routing uncertainty. We model the routing decision as a multinomial choice over the $K$ experts. Let the prior over the expert ensembling coefficients $\alpha$ be modeled as a symmetric Dirichlet distribution:
    \begin{align}
        p(\alpha \mid \gamma_0) &= \text{Dir}(\alpha \mid \gamma_0 \mathbf{1}_K) \nonumber \\
        &= \frac{\Gamma(K \gamma_0)}{\Gamma(\gamma_0)^K} \prod_{k=1}^K \alpha_k^{\gamma_0 - 1}
    \end{align}
    where $\gamma_0 > 0$ represents the prior concentration representing baseline confidence. During inference, we observe the task-specific environmental attraction coordinates $u_{b} = [u_{1,b}, \dots, u_{K,b}]^T \in [0, 1]^K$. We interpret these normalized affinities as empirical pseudocount observations $\mathbf{n}_b = [n_{1,b}, \dots, n_{K,b}]^T$, where $n_{k,b} = \kappa \cdot u_{k,b}$, and $\kappa > 0$ represents the observation scaling intensity (precision). By Bayes' rule, the posterior distribution over the routing coefficients is conjugate and also Dirichlet:
    \begin{equation}
        p(\alpha \mid \mathbf{n}_b, \gamma_0) = \text{Dir}(\alpha \mid \beta_{k,b})
    \end{equation}
    where the posterior concentration parameters are given by $\beta_{k,b} = \gamma_0 + \kappa \cdot u_{k,b}$.
    
    The total concentration of the posterior Dirichlet distribution is $S_b = \sum_{k=1}^K \beta_{k,b} = K \gamma_0 + \kappa \sum_{k=1}^K u_{k,b}$. In Bayesian statistics, the total concentration $S_b$ directly represents the total evidence or information content of the observations. We can define the \emph{Bayesian Confidence} $C^{\text{Bayes}}_b \in [0, 1]$ as the proportion of total concentration contributed by the empirical observations relative to the non-informative prior:
    \begin{equation}
        C^{\text{Bayes}}_b = \frac{S_b - K \gamma_0}{S_b} = \frac{\kappa \sum_{k=1}^K u_{k,b}}{K \gamma_0 + \kappa \sum_{k=1}^K u_{k,b}}
    \end{equation}
    When the task affinities are low or inconsistent (high OOD noise), $\sum_k u_{k,b} \to 0$, leading to $C^{\text{Bayes}}_b \to 0$. Conversely, when there is strong, consistent task-specific evidence, $\sum_k u_{k,b}$ is large, leading to $C^{\text{Bayes}}_b \to 1$.
    
    We can then formulate the dynamic sharpening strength $\gamma_{\text{dais}, b}$ directly using this self-calibrated Bayesian Confidence:
    \begin{equation}
        \gamma_{\text{dais}, b} = 1.0 + (\gamma_{\max} - 1.0) \cdot C^{\text{Bayes}}_b
    \end{equation}
    This formulation is completely parameter-free with respect to the exponential decay rate $\eta_{\text{decay}}$, as the confidence scaling is driven entirely by the relative Bayesian posterior concentration of task affinities. By substituting $C^{\text{Bayes}}_b$ into our expected utility formulation, we achieve an elegant, mathematically airtight, and self-calibrating alternative that resolves any empirical hyperparameter tuning critiques, providing a robust probabilistic grounding for dynamic ensembling. We view this fully probabilistic, self-calibrating routing formulation as an exceptionally promising and theoretically satisfying avenue for future research.
\end{enumerate}

By decoupling the continuous dynamics of the solver from the prediction inference stage, DAIS and its adaptive E-ITAS extension allow the model to benefit from robust, noise-robust lateral cooperation and self-regulating competitive exclusion during routing, while executing sharp, dilution-free classification decisions.

By utilizing only basic tensor additions, multiplications, and a small matrix multiplication, DESS incurs negligible latency (fewer than 5 milliseconds) and requires no backpropagation or parameter updates.

We clarify that the Lotka-Volterra dynamics act as an \emph{internal} recurrent feedback loop over the blending coefficients $\alpha^{(t)}$, rather than a closed-loop feedback from downstream layers. In other words, the environmental attraction vector $u_{k, b}$ is computed once at the routing layer and held constant, while the ensembling coefficients dynamically self-organize through the competitive-cooperative interactions of the solver. This localized recurrent design ensures that the solver converges rapidly inside a single forward pass without requiring multiple expensive forward and backward passes of downstream transformer blocks, which would introduce prohibitive serving latency.

\subsection{Paradox-Free Execution Layout and Activation Blending}
\label{subsec:architecture}
We deploy ESM-LVC in a \textbf{Paradox-Free Execution Layout} designed to resolve representation bottlenecks in early layers. We divide a standard Vision Transformer (ViT) consisting of $L$ blocks into shared and specialized regions:
\begin{enumerate}
    \item \textbf{Shared Feature Extractor (Layers 1 to $L_{\text{route}}$):} The first $L_{\text{route}}$ blocks are frozen and shared across all experts during training, containing zero LoRA adapters. This allows the model to extract robust, task-invariant representation centroids.
    \item \textbf{Routing Layer ($L_{\text{route}}$):} Located at block $L_{\text{route}}$ (e.g., Layer 3). Here, the intermediate pooled features $h^{(L_{\text{route}})}_b$ are extracted to compute $u_{k, b}$ and run DESS to solve the final coefficients $\alpha^{\text{final}}_{k, b}$.
    \item \textbf{Specialized Expert Region (Layers $L_{\text{route}} + 1$ to $L$):} These layers contain task-specific LoRA adapters. Inside these blocks, the adapter activations are dynamically blended sample-wise using $\alpha^{\text{final}}_{k, b}$ in a single, parallel forward pass:
    \begin{equation}
        h_b^{(l)} = h_b^{(l-1)} W_{\text{base}}^{(l)} + \text{DSA}\left( \sum_{k=1}^K \alpha^{\text{final}}_{k, b} y_{k, b}^{(l)} \right)
    \end{equation}
    where $y_{k, b}^{(l)} = h_b^{(l-1)} A_k^{(l)} B_k^{(l)}$ represents the individual activation output of the LoRA adapter for expert $k$ at layer $l$, $A_k^{(l)}$ and $B_k^{(l)}$ are its down- and up-projection matrices, and $\text{DSA}(\cdot)$ is the Dynamic Scale Alignment operator described below.
\end{enumerate}

\paragraph{Dynamic Scale Alignment (DSA) for Activation Scale Stability:}
In physical transformers, linearly combining activations from multiple unaligned LoRA adapters can distort vector magnitudes and trigger scale drift, violating the statistics expected by downstream LayerNorm/RMSNorm layers. To ensure absolute scale stability, we introduce a \textbf{Dynamic Scale Alignment (DSA)} operator. Let $y_{k, b}^{(l)}$ represent the individual activation output of LoRA expert $k$ for sample $b$. The blended adapter activation is given by $y^{\text{blend}}_b = \sum_k \alpha^{\text{final}}_{k, b} y_{k, b}^{(l)}$. We formulate DSA as:
\begin{equation}
    \text{DSA}\left( y^{\text{blend}}_b \right) = y^{\text{blend}}_b \cdot \frac{\sum_k \alpha^{\text{final}}_{k, b} \| y_{k, b}^{(l)} \|_2}{\| y^{\text{blend}}_b \|_2 + \epsilon}
\end{equation}
where $\epsilon = 10^{-8}$ is a numerical stabilizer. This dynamic rescaling operator guarantees that the norm of the blended activation matches the expected norm of a single active adapter, completely preventing scale dampening and out-of-distribution magnitude drift across downstream blocks, preserving the model's structural integrity.

Finally, the blended features are passed through the expert-specific classification heads to output predicted labels. This layout ensures that we achieve single-pass, dynamic multi-adapter serving without routing latency or structural conflicts.
